% Esquisse d'un programme (cote 119, copy 1): archive cote 119, site page 18 (PDF page 19), typescript page 16.
% typescript. Draft compared: EsquisseFr.pdf (text layer):601-668.
% AI-transcribed by Claude Opus 5.5 (claude-opus-5-5, adaptive thinking, effort high), Message Batches API, batch esq1 (msgbatch_01LNsogFyF3y6KSp9UbtJtBk), 2026-11-12; not yet checked by eye.
\scanpage{119/18}
\mspage{16}

qu'apparaît le thème central de la géométrie algébrique anabélienne,
qui est de reconstituer certaines variétés $X$ dites ``anabéliennes''
sur un corps absolu $K$ à partir de leur groupe fondamental mixte,
extension de $\mathrm{Gal}(\bar{K}/K)$ par $\pi_1(X_{\bar{K}})$ ; c'est alors que se dégage la
``conjecture fondamentale de la géométrie algébrique anabélienne'',
proche des conjectures de Mordell et de Tate que vient de démontrer
Faltings($^3$). C'est là aussi que s'amorcent une première réflexion sur
les groupes de Teichmüller, et les premières intuitions sur la struc-%
ture multiple de la ``tour de Teichmüller'' - les multiplicités modu-%
laires ouvertes $M_{g,\nu}$ apparaissant par ailleurs comme les premiers
exemples importants, en dimension $> 1$, de variétés (ou plutôt, de
multiplicités) qui semblent bien mériter l'appellation ``anabélienne''.
Vers la fin de cette période de réflexion, celle-ci m'apparaît com-%
me une réflexion fondamentale sur une théorie alors encore dans les
limbes, pour laquelle l'appellation ``\uline{théorie de Galois-Teichmüller}''
me semble plus appropriée que ``Théorie de Galois'' que j'avais d'abord
donnée à mes notes.

Ce n'est pas le lieu ici de donner un aperçu plus circons-%
tancié de cet ensemble de questions, intuitions, idées - y compris
des résultats palpables, certes. Le plus important me semble celui
signalé en passant au par. 2, savoir la fidélité de l'action exté-%
rieure de $\mathrm{I\Gamma} = \mathrm{Gal}(\bar{\mathbb{Q}}/\mathbb{Q})$ (et de ses sous-groupes ouverts) sur $\hat{\pi}_{0,3}$,
et plus généralement (si je me rappelle bien) sur le groupe fonda-%
mental de toute courbe algébrique ``anabélienne'' (i.e. dont le genre
$g$ et le ``nombre de trous'' $\nu$ satisfont l'inégalité $2g + \nu \geq 3$,
i.e. telle que $\chi(X) < 0$) définie sur une extension finie de $\mathbb{Q}$. Ce
résultat peut être considéré comme essentiellement équivalent au
théorème de Bielyi - c'est la première manifestation concrète, par
un énoncé mathématique précis, du ``message'' dont il a été question
plus haut.

Je voudrais terminer cet aperçu rapide par quelques mots
de commentaire sur la richesse vraiment inimaginable d'un groupe
anabélien typique comme le groupe $S\ell(2,\mathbb{Z})$ - sans doute le groupe
discret infini le plus remarquable qu'on ait rencontré, qui apparaît

\margin{18}
